\chapter*{Удиртгал}
\addcontentsline{toc}{chapter}{УДИРТГАЛ}
Орчин үед байгууллага, компаниуд хэрэглэгчдийнхээ санал сэтгэгдлийг сонсох, тренд болж буй сэдвүүдийг илрүүлэх зорилгоор сошиал мониторингийн системийг өргөнөөр ашиглах болсон. Одоогийн байдлаар Sumbee.mn зэрэг системүүд нь ихэвчлэн түлхүүр үг дээр тулгуурласан тоон болон эерэг/сөрөг (sentiment) анализ хийдэг бөгөөд агуулгын түвшинд задлан шинжлэх, илүү бодитой, бүтэцчилсэн дүгнэлт гаргах боломж дутмаг байна. Мөн маш их хэмжээний өгөгдөл дээр LLM шууд ашиглах нь тооцооллын зардал маш их шаардах ба удаан (latency өндөр) байдаг. Иймд Монгол хэл, тэр дундаа сошиал өгөгдөл дээр үр дүнтэй ажиллаж чадах NLP (Байгалийн хэлний боловсруулалт)-д суурилсан шийдэл боловсруулах шаардлага тулгарч байна.

Энэхүү судалгааны ажлын зорилго нь сошиал сүлжээний бүтэцлэгдээгүй текст өгөгдлөөс далд сэдвүүд (Topic) болон нэрлэсэн объектуудыг (NER) автоматаар илрүүлж, тэдгээрийн хамаарлын сүлжээг (Network Analysis) үүсгэх инференс пайплайн зохиож хэрэгжүүлэхэд оршино.

\newpage
\chapter*{Нэр томьёоны тайлбар}
\addcontentsline{toc}{chapter}{Нэр томьёоны тайлбар}
\begin{itemize}
    \item \textbf{Topic Modeling}: Сэдэв загварчлал
    \item \textbf{NLP}: Байгалийн хэлний боловсруулалт (Natural Language Processing)
    \item \textbf{Text Mining}: Текст олборлолт
    \item \textbf{Embedding}: Векторчилсон хэсгийн төлөөлөл
\end{itemize}

\newpage
\chapter*{Товчилсон үгс}
\addcontentsline{toc}{chapter}{Товчилсон үгс}
\begin{itemize}
    \item \textbf{LLM}: Large Language Model (Том хэлний загвар)
    \item \textbf{NER}: Named Entity Recognition 
    \item \textbf{API}: Application Programming Interface
    \item \textbf{EDA}: Event-Driven Architecture
    \item \textbf{DB}: Database (Өгөгдлийн сан)
\end{itemize}
